%===================== REFERENCIAL TEÓRICO =====================
\section{Referencial teórico}

% ESTA É A SEÇÃO QUE O ORIENTADOR MAIS COBROU.
% Ela se escreve a partir dos (até) 10 artigos selecionados na busca.
% Para CADA artigo, anotar: o que mediram, com quantas pessoas, que
% resultado deu. Depois agrupar por tema nas subseções abaixo.

\subsection{Fobias específicas a animais}

% O que são, prevalência, critérios. Base conceitual.
[ESCREVER]

\subsection{Exposição gradual}

% O princípio: contato progressivo, controlado, com o estímulo temido.
% Por que a graduação importa e por que o susto é contraproducente.
[ESCREVER]

\subsection{Realidade Virtual aplicada à exposição}

% O núcleo da revisão. Aqui entram os estudos encontrados: o que já foi
% feito, com que equipamento, com que resultado.
[ESCREVER]

\subsection{Realidade Virtual no navegador}

% WebXR, A-Frame, execução sem instalação. É aqui que se argumenta a
% contribuição deste trabalho frente ao que a literatura já cobriu.
[ESCREVER]

\subsection{Síntese dos trabalhos relacionados}

% Fechar com uma tabela comparativa: trabalho | estímulo | equipamento |
% n | o que mediram. É o que a banca mais olha.
[ESCREVER]
